\documentclass[a4paper,12pt]{article}
\usepackage[T2A]{fontenc}
\usepackage[utf8]{inputenc}
\usepackage[russian]{babel}
\usepackage{amsmath,amssymb,mathtools}
\renewcommand{\familydefault}{\sfdefault}
\usepackage{icomma}
\usepackage[a4paper,left=22mm,right=22mm,top=23mm,bottom=22mm,headheight=15pt]{geometry}
\usepackage{microtype}
\usepackage{xcolor}
\usepackage{tikz}
\usetikzlibrary{calc,arrows.meta,positioning,shapes.geometric}
\usepackage{pgfplots}
\pgfplotsset{compat=1.18}
\usepackage{enumitem,tabularx,booktabs}
\usepackage{fancyhdr}
\usepackage{setspace}
\usepackage{hyperref}
\hypersetup{hidelinks,pdftitle={Чек-лист. Степени, прикладные задачи и показательные функции},pdfauthor={Лёвин Артём Александрович}}
\definecolor{accent}{HTML}{137A7A}
\definecolor{darkaccent}{HTML}{0A525A}
\definecolor{textgray}{HTML}{3F4A4D}
\definecolor{pale}{HTML}{E8EFEF}
\color{textgray}
\onehalfspacing
\setlength{\parindent}{0pt}
\setlength{\parskip}{5pt}
\setcounter{secnumdepth}{0}
\setlist[enumerate]{leftmargin=1.75em,itemsep=4pt,topsep=5pt}
\pagestyle{fancy}
\fancyhf{}
\fancyhead[L]{\small\color{darkaccent}Марк Гукин \quad\textbar\quad 11.10.2026}
\fancyhead[R]{\small\color{textgray}Степени и функции}
\fancyfoot[C]{\small\thepage}
\renewcommand{\headrulewidth}{0.35pt}
\renewcommand{\footrulewidth}{0pt}
\newcommand{\sectiontitle}[1]{\vspace{4mm}{\Large\bfseries\color{darkaccent}#1}\par\vspace{2mm}\color{textgray}}
\newcommand{\subtopic}[1]{\vspace{1mm}{\large\bfseries\color{darkaccent}#1}\par\color{textgray}}
\newcommand{\key}[1]{\textbf{\color{darkaccent}#1}}
\newcommand{\goodcheck}{\ensuremath{\square}}
\begin{document}

% ТИТУЛ
\thispagestyle{empty}
\begin{center}
\vspace*{33mm}
{\small\bfseries\color{darkaccent} ИНДИВИДУАЛЬНОЕ УЧЕБНОЕ ПОСОБИЕ}\par
\vspace{18mm}
{\fontsize{35.83}{42}\selectfont\bfseries\color{darkaccent}Чек-лист}\par
\vspace{8mm}
{\Large\bfseries Степени и показательные функции}\par
\vspace{3mm}
{\large Прикладные задачи. Графики. Вычисления.}\par
\vspace{22mm}
{\large Марк Гукин}\par
\vspace{7mm}
{\normalsize 11 октября 2026 года}\par
\vfill
{\normalsize\color{textgray}Лёвин Артём Александрович}\par
{\normalsize\bfseries\color{darkaccent}эксклюзивно для Марка Гукина}\par
\vspace{17mm}
\begin{tikzpicture}[x=1cm,y=1cm]
\draw[darkaccent,opacity=.48,line width=.95pt]
(-7.4,0) .. controls (-4.6,.36) and (-3.1,-.38) .. (-.4,.13)
.. controls (2.2,.55) and (4.4,-.36) .. (7.4,.04);
\end{tikzpicture}
\vspace{7mm}
\end{center}
\clearpage

\sectiontitle{01 \quad Степени: опорные правила}
Для любого $a\ne0$ и целых $m,n$ действуют свойства:
\[
\begin{aligned}
 a^m a^n &=a^{m+n},&\qquad \frac{a^m}{a^n}&=a^{m-n},\\[2pt]
 (a^m)^n&=a^{mn},&\qquad a^{-n}&=\frac1{a^n},\\[2pt]
 (ab)^n&=a^n b^n,&\qquad a^0&=1.
\end{aligned}
\]
\key{Проверка.} При делении степеней показатели \emph{вычитаются}; при возведении степени в степень --- \emph{перемножаются}. Отрицательная степень означает обратное число, а не отрицательный знак перед основанием.

\subtopic{Пример. Приведение степеней к общему основанию}
\[
\frac{10^{10}}{2^6\cdot5^8}
=\frac{2^{10}5^{10}}{2^6 5^8}=2^4 5^2=400.
\]
Удобный приём для больших чисел: $10^n=(2\cdot5)^n=2^n5^n$.

\sectiontitle{02 \quad Прикладная задача: газ}
В задаче использована зависимость давления от объёма
\[
\boxed{pV^a=\mathrm{const}},\qquad p>0,\quad V>0.
\]
Здесь $p$ --- давление, $V$ --- объём, $a$ --- неизвестный показатель. Для двух состояний одного процесса
\[
\boxed{p_1V_1^a=p_2V_2^a}.
\]
\key{Пример.} При уменьшении объёма вдвое давление увеличилось вчетверо. Найдите $a$.

Поскольку $V_1=2V_2$ и $p_2=4p_1$, получаем
\[
 p_1(2V_2)^a=4p_1V_2^a
 \quad\Longrightarrow\quad 2^a=4=2^2
 \quad\Longrightarrow\quad\boxed{a=2}.
\]
\key{Самопроверка.} Множители $p_1$ и $V_2^a$ можно сократить: они ненулевые.

\clearpage
\sectiontitle{03 \quad Прикладная задача: температура звезды}
Вторая рассмотренная модель описывается формулой
\[
\boxed{P=\sigma ST^4},
\]
где $P$ --- излучаемая мощность, $S$ --- площадь поверхности, $T$ --- абсолютная температура, $\sigma$ --- заданная постоянная.

\key{Пример из занятия.} Дано:
\[
\sigma=5{,}7\cdot10^{-8},\qquad
S=\frac{10^{20}}{228},\qquad
P=1{,}5625\cdot10^{25}.
\]
Найдём температуру $T>0$:
\[
T^4=\frac{P}{\sigma S}
=\frac{1{,}5625\cdot10^{25}\cdot228}
       {5{,}7\cdot10^{-8}\cdot10^{20}}
=6{,}25\cdot10^{14}.
\]
Представим число как четвёртую степень:
\[
6{,}25\cdot10^{14}
=625\cdot10^{12}=5^4\cdot(10^3)^4=(5\cdot10^3)^4.
\]
Итак,
\[
\boxed{T=5000\ \text{К}}.
\]
\key{Контроль.} При извлечении корня четвёртой степени из $T^4$ учитываем физическое условие $T>0$.

\vspace{3mm}
\subtopic{Три полезных приёма}
\begin{enumerate}[label=\textbf{\arabic*.}]
\item Выразить искомую величину: $T^4=P/(\sigma S)$.
\item Собрать множители и степени десятки до вычисления больших чисел.
\item Распознать полную четвёртую степень и проверить знак результата.
\end{enumerate}

\vspace{4mm}
\sectiontitle{04 \quad Показательная функция}
Рабочая модель для графиков:
\[
\boxed{f(x)=a^x+b},\qquad a>0,\quad a\ne1.
\]
При $x=0$ имеем $a^0=1$, поэтому \key{$b=f(0)-1$}. Значение $b$ также задаёт горизонтальную асимптоту $y=b$.
\clearpage

\sectiontitle{05 \quad Восстанавливаем функцию по точкам}
\key{Метод узловых точек:} выбираем две точки графика с известными целыми координатами; особенно удобна точка с $x=0$.

\begin{minipage}[t]{.49\linewidth}
\vspace{0pt}
\textbf{Случай 1. Основание $a>1$}\par
\smallskip
\begin{tikzpicture}
\begin{axis}[
width=\linewidth,height=61mm,scale only axis=false,
xmin=-2.4,xmax=2.7,ymin=-3.8,ymax=4.4,
axis lines=middle,xtick={-2,-1,0,1,2},ytick={-3,-2,-1,0,1,2,3,4},
minor tick num=0,grid=major,grid style={line width=.25pt,draw=gray!25},
axis line style={black!70},ticklabel style={font=\scriptsize},
clip=false]
\addplot[accent,very thick,domain=-2.3:2.5,samples=170]{2^x-3};
\addplot[only marks,mark=*,mark size=2.5pt,color=darkaccent] coordinates {(0,-2)(1,-1)};
\draw[gray,dashed] (axis cs:-2.3,-3)--(axis cs:2.65,-3);
\node[font=\scriptsize,anchor=north east,fill=white,inner sep=1pt] at (axis cs:0,-2){$A(0;-2)$};
\node[font=\scriptsize,anchor=north west,fill=white,inner sep=1pt] at (axis cs:1,-1){$B(1;-1)$};
\end{axis}
\end{tikzpicture}
\end{minipage}\hfill
\begin{minipage}[t]{.49\linewidth}
\vspace{0pt}
\textbf{Случай 2. Основание $0<a<1$}\par
\smallskip
\begin{tikzpicture}
\begin{axis}[
width=\linewidth,height=61mm,scale only axis=false,
xmin=-3.0,xmax=2.5,ymin=-4.7,ymax=5.3,
axis lines=middle,xtick={-3,-2,-1,0,1,2},ytick={-4,-2,0,2,4},
minor tick num=0,grid=major,grid style={line width=.25pt,draw=gray!25},
axis line style={black!70},ticklabel style={font=\scriptsize},
clip=false]
\addplot[accent,very thick,domain=-3.0:2.5,samples=170]{2^(-x)-4};
\addplot[only marks,mark=*,mark size=2.5pt,color=darkaccent] coordinates {(0,-3)(-2,0)};
\draw[gray,dashed] (axis cs:-3,-4)--(axis cs:2.45,-4);
\node[font=\scriptsize,anchor=north west,fill=white,inner sep=1pt] at (axis cs:0,-3){$A(0;-3)$};
\node[font=\scriptsize,anchor=south west,fill=white,inner sep=1pt] at (axis cs:-2,0){$B(-2;0)$};
\end{axis}
\end{tikzpicture}
\end{minipage}

\vspace{3mm}
\begin{minipage}[t]{.49\linewidth}
\textbf{Решение для левого графика}\par
\[
\begin{cases}
-2=a^0+b,\\[-1mm]
-1=a^1+b.
\end{cases}
\Rightarrow\begin{cases}b=-3,\\a=2.
\end{cases}
\]
\[
 f(x)=2^x-3,\qquad
 \boxed{f(6)=64-3=61}.
\]
\end{minipage}\hfill
\begin{minipage}[t]{.49\linewidth}
\textbf{Решение для правого графика}\par
\[
\begin{cases}
-3=a^0+b,\\[-1mm]
0=a^{-2}+b.
\end{cases}
\Rightarrow\begin{cases}b=-4,\\a^{-2}=4.
\end{cases}
\]
\[
 a=\frac12,\quad f(x)=2^{-x}-4.
\]
\end{minipage}

\vspace{3mm}
\key{Продолжение.} Если $f(x)=2^{-x}-4$ и $f(x)=12$, то
\[
2^{-x}-4=12\quad\Rightarrow\quad
2^{-x}=16=2^4\quad\Rightarrow\quad\boxed{x=-4}.
\]
\key{Проверяйте подстановкой:} $f(-4)=2^4-4=12$. Показатель $-x$ требует особого внимания к знаку.

\clearpage
\sectiontitle{Алгоритм: функция по двум узловым точкам}
\vspace{12mm}
\begin{center}
\begin{tikzpicture}[
node distance=10mm and 12mm,
endnode/.style={ellipse,draw=darkaccent,line width=.85pt,align=center,minimum width=37mm,minimum height=12mm,fill=pale},
action/.style={rounded corners=2pt,draw=darkaccent,line width=.85pt,align=center,text width=113mm,minimum height=16mm,inner sep=5pt},
arrow/.style={-{Stealth[length=3mm]},line width=.9pt,draw=darkaccent}
]
\node[endnode] (start) {Начало};
\node[action,below=of start] (n1) {Запишите вид функции $f(x)=a^x+b$\\ $(a>0,\ a\ne1)$};
\node[action,below=of n1] (n2) {Найдите на графике две узловые точки\\ $A(0;y_0)$ и $B(k;y_1)$, где $k\ne0$};
\node[action,below=of n2] (n3) {По точке $A$ найдите сдвиг: $b=y_0-1$};
\node[action,below=of n3] (n4) {По точке $B$ найдите основание: $a^k=y_1-b$};
\node[action,below=of n4] (n5) {Проверьте $a>0$, $a\ne1$ и подставьте обе точки};
\node[action,below=of n5] (n6) {Вычислите требуемое значение $f(t)$\\ или решите уравнение $f(x)=c$};
\node[endnode,below=of n6] (stop) {Ответ};
\foreach \a/\b in {start/n1,n1/n2,n2/n3,n3/n4,n4/n5,n5/n6,n6/stop}{\draw[arrow] (\a)--(\b);}
\end{tikzpicture}
\end{center}
\clearpage

\sectiontitle{Тренировочный блок \quad \textnormal{Путь А}}
\textit{10 заданий. Выполняйте по порядку, записывайте основные преобразования и ответ.}
\vspace{3mm}
\begin{enumerate}[label=\textbf{\arabic*.},itemsep=10pt,parsep=0pt]
\item Вычислите: $\displaystyle\frac{2^8\cdot 2^{-3}}{2^2}$.
\item Вычислите: $\displaystyle\frac{(3^2)^3}{3^4}$.
\item Вычислите: $\displaystyle\left(\frac{2^{-1}}{2^2}\right)^{-2}$.
\item В процессе $pV^a=\mathrm{const}$ объём уменьшился в $2$ раза, а давление возросло в $8$ раз. Найдите $a$.
\item В процессе $pV^a=\mathrm{const}$ объём увеличился в $4$ раза, а давление уменьшилось в $16$ раз. Найдите $a$.
\item График $f(x)=a^x+b$ проходит через точки $A(0;-2)$ и $B(1;1)$. Найдите $a$, $b$ и $f(2)$.
\item График $f(x)=a^x+b$ проходит через точки $A(0;-1)$ и $B(-2;2)$. Найдите $a$, $b$ и $f(2)$.
\item Дана функция $f(x)=2^x-5$. Найдите $x$, если $f(x)=11$.
\item Дана функция $f(x)=3^{-x}-2$. Найдите $x$, если $f(x)=25$.
\item Излучение описывается формулой $P=\sigma ST^4$. Найдите $T>0$, если
\[
\sigma=5\cdot10^{-8},\quad S=2\cdot10^3,\quad P=10^8.
\]
\end{enumerate}
\vspace{5mm}
{\small\color{darkaccent}\textbf{Перед сдачей:} \quad
$\square$ проверены правила степеней \quad
$\square$ учтены знаки показателей \quad
$\square$ проверены подстановки.}
\clearpage

\sectiontitle{Ответы к тренировочному блоку}
\vspace{6mm}
\renewcommand{\arraystretch}{1.6}
\begin{tabularx}{\linewidth}{@{}cX@{}}
\toprule
\textbf{№} & \textbf{Ответ} \\
\midrule
1 & $8$ \\
2 & $9$ \\
3 & $64$ \\
4 & $a=3$ \\
5 & $a=2$ \\
6 & $a=4$, $b=-3$, $f(2)=13$ \\
7 & $a=\frac12$, $b=-2$, $f(2)=-\frac74$ \\
8 & $x=4$ \\
9 & $x=-3$ \\
10 & $T=1000\ \text{К}$ \\
\bottomrule
\end{tabularx}
\vspace{15mm}
\subtopic{Быстрая самопроверка}
\begin{enumerate}[label=\textbf{\arabic*.},itemsep=8pt]
\item Подставьте найденные $a,b$ в координаты обеих точек графика.
\item В задачах на газ сравните изменение $p$ и $V^a$.
\item Для температуры проверьте равенство $P=\sigma ST^4$.
\end{enumerate}
\end{document}
